% GENERATED FILE. Edit canonical lessons, then run npm run generate:books.
% canonical-source-hash: fnv1a64:c3d98667aa19758d
% canonical-lessons: GU-C473-jagh, GU-C473-taliyu, GU-C473-duti, GU-C473-pindi, GU-C473-bagal

\chapter{From Hip to Toe}
\label{ch:gu-from-hip-to-toe}
\hlchaptermodality{473}

\begin{chapteropening}
\textbf{By the end of this chapter:} \emph{I can name parts of the leg and arm.}

Complete the last lesson of chapter 473: I can name parts of the leg and arm.
\end{chapteropening}

\section[\texorpdfstring{\gu{જાંઘ} (jā̃gh)}{જાંઘ (jā̃gh)}]{\gu{જાંઘ} (jā̃gh) — a thigh}

\label{lesson:GU-C473-jagh}



\begin{quote}
Before the new one: say the Gujarati for an intestine, then the Gujarati for the stomach.
\end{quote}

\subsection*{You'll want to know: \gu{જાંઘ}}
\textbf{\gu{જાંઘ}} — \emph{jā̃gh} — \textquotedblleft{}a thigh\textquotedblright{}.

\begin{quote}
\gu{પડ્યા} \gu{પછી} \gu{મારી} \gu{જાંઘ} \gu{પર} \gu{સોજો} \gu{આવ્યો}.
\end{quote}

\begin{quote}
\emph{paḍyā pachī mārī jā̃gh par sojo āvyo.}
\end{quote}

\begin{quote}
After I fell, my thigh swelled up.
\end{quote}

A feminine noun.

\begin{cousinweb}
\textbf{\gu{જાંઘ}} (\emph{jā̃gh}) is inherited from Sanskrit \emph{jaṅghā}, \textquotedblleft{}shank, ankle to knee\textquotedblright{}; CDIAL lists the Gujarati word, \textquotedblleft{}thigh\textquotedblright{} (CDIAL 5082).
\end{cousinweb}

\begin{grammarlens}[title={What you've built}]
One more word for describing what you lived and felt.
\end{grammarlens}

\subsection*{Your turn}
\begin{itemize}
\raggedright
  \item \emph{Say it:} \emph{jā̃gh}
  \item \emph{Say it:} \emph{jā̃gh}, once more
  \item \emph{Say it:} \emph{paḍyā pachī mārī jā̃gh par sojo āvyo.} — After I fell, my thigh swelled up.
  \item \emph{Recall:} say the Gujarati for an intestine, then the Gujarati for the stomach, then say \emph{jā̃gh} again
\end{itemize}

\begin{tcolorbox}[breakable,colback=teal!4,colframe=teal!35!black,title={Before you move on}]
What does \gu{જાંઘ} mean? (\textquotedblleft{}A thigh\textquotedblright{}.) Say it once more.
\end{tcolorbox}
\section[\texorpdfstring{\gu{તળિયું} (taḷiyũ)}{તળિયું (taḷiyũ)}]{\gu{તળિયું} (taḷiyũ) — the sole of the foot}

\label{lesson:GU-C473-taliyu}



\begin{quote}
Before the new one: say the Gujarati for the stomach, then the Gujarati for a thigh.
\end{quote}

\subsection*{You'll want to know: \gu{તળિયું}}
\textbf{\gu{તળિયું}} — \emph{taḷiyũ} — \textquotedblleft{}the sole of the foot\textquotedblright{}.

\begin{quote}
\gu{ગરમ} \gu{રેતીથી} \gu{પગનાં} \gu{તળિયાં} \gu{બળ્યાં}.
\end{quote}

\begin{quote}
\emph{garam retīthī pagnā̃ taḷiyā̃ baḷyā̃.}
\end{quote}

\begin{quote}
The hot sand burned the soles of my feet.
\end{quote}

A neuter noun.

\begin{cousinweb}
\textbf{\gu{તળિયું}} (\emph{taḷiyũ}) is inherited from Sanskrit \emph{tala}, \textquotedblleft{}base, bottom\textquotedblright{}; CDIAL lists the Gujarati word, \textquotedblleft{}bottom, flat surface, sole\textquotedblright{} (CDIAL 5731).
\end{cousinweb}

\begin{grammarlens}[title={What you've built}]
One more word for describing what you lived and felt.
\end{grammarlens}

\subsection*{Your turn}
\begin{itemize}
\raggedright
  \item \emph{Say it:} \emph{taḷiyũ}
  \item \emph{Say it:} \emph{taḷiyũ}, once more
  \item \emph{Say it:} \emph{garam retīthī pagnā̃ taḷiyā̃ baḷyā̃.} — The hot sand burned the soles of my feet.
  \item \emph{Recall:} say the Gujarati for the stomach, then the Gujarati for a thigh, then say \emph{taḷiyũ} again
\end{itemize}

\begin{tcolorbox}[breakable,colback=teal!4,colframe=teal!35!black,title={Before you move on}]
What does \gu{તળિયું} mean? (\textquotedblleft{}The sole of the foot\textquotedblright{}.) Say it once more.
\end{tcolorbox}
\section[\texorpdfstring{\gu{ડૂંટી} (ḍū̃ṭī)}{ડૂંટી (ḍū̃ṭī)}]{\gu{ડૂંટી} (ḍū̃ṭī) — the navel}

\label{lesson:GU-C473-duti}



\begin{quote}
Before the new one: say the Gujarati for a thigh, then the Gujarati for the sole of the foot.
\end{quote}

\subsection*{You'll want to know: \gu{ડૂંટી}}
\textbf{\gu{ડૂંટી}} — \emph{ḍū̃ṭī} — \textquotedblleft{}the navel\textquotedblright{}.

\begin{quote}
\gu{બાળકની} \gu{ડૂંટી} \gu{પર} \gu{દાક્તરે} \gu{મલમ} \gu{લગાડ્યો}.
\end{quote}

\begin{quote}
\emph{bāḷaknī ḍū̃ṭī par dāktare malam lagāḍyo.}
\end{quote}

\begin{quote}
The doctor put ointment on the baby's navel.
\end{quote}

A feminine noun.

\begin{cousinweb}
\textbf{\gu{ડૂંટી}} (\emph{ḍū̃ṭī}) is inherited from Sanskrit \emph{tundi}, \textquotedblleft{}navel\textquotedblright{}; CDIAL lists the Gujarati word (CDIAL 5860).
\end{cousinweb}

\begin{grammarlens}[title={What you've built}]
One more word for describing what you lived and felt.
\end{grammarlens}

\subsection*{Your turn}
\begin{itemize}
\raggedright
  \item \emph{Say it:} \emph{ḍū̃ṭī}
  \item \emph{Say it:} \emph{ḍū̃ṭī}, once more
  \item \emph{Say it:} \emph{bāḷaknī ḍū̃ṭī par dāktare malam lagāḍyo.} — The doctor put ointment on the baby's navel.
  \item \emph{Recall:} say the Gujarati for a thigh, then the Gujarati for the sole of the foot, then say \emph{ḍū̃ṭī} again
\end{itemize}

\begin{tcolorbox}[breakable,colback=teal!4,colframe=teal!35!black,title={Before you move on}]
What does \gu{ડૂંટી} mean? (\textquotedblleft{}The navel\textquotedblright{}.) Say it once more.
\end{tcolorbox}
\section[\texorpdfstring{\gu{પિંડી} (pinḍī)}{પિંડી (pinḍī)}]{\gu{પિંડી} (pinḍī) — the calf of the leg}

\label{lesson:GU-C473-pindi}



\begin{quote}
Before the new one: say the Gujarati for the sole of the foot, then the Gujarati for the navel.
\end{quote}

\subsection*{You'll want to know: \gu{પિંડી}}
\textbf{\gu{પિંડી}} — \emph{pinḍī} — \textquotedblleft{}the calf of the leg\textquotedblright{}.

\begin{quote}
\gu{લાંબું} \gu{ચાલ્યા} \gu{પછી} \gu{મારી} \gu{પિંડી} \gu{દુખે} \gu{છે}.
\end{quote}

\begin{quote}
\emph{lā̃bũ cālyā pachī mārī pinḍī dukhe che.}
\end{quote}

\begin{quote}
After a long walk my calf aches.
\end{quote}

A feminine noun.

\begin{cousinweb}
\textbf{\gu{પિંડી}} (\emph{pinḍī}) is inherited from Sanskrit \emph{piṇḍa}, \textquotedblleft{}lump\textquotedblright{}, which also meant \textquotedblleft{}calf of the leg\textquotedblright{}; CDIAL lists the Gujarati word (CDIAL 8168).
\end{cousinweb}

\begin{grammarlens}[title={What you've built}]
One more word for describing what you lived and felt.
\end{grammarlens}

\subsection*{Your turn}
\begin{itemize}
\raggedright
  \item \emph{Say it:} \emph{pinḍī}
  \item \emph{Say it:} \emph{pinḍī}, once more
  \item \emph{Say it:} \emph{lā̃bũ cālyā pachī mārī pinḍī dukhe che.} — After a long walk my calf aches.
  \item \emph{Recall:} say the Gujarati for the sole of the foot, then the Gujarati for the navel, then say \emph{pinḍī} again
\end{itemize}

\begin{tcolorbox}[breakable,colback=teal!4,colframe=teal!35!black,title={Before you move on}]
What does \gu{પિંડી} mean? (\textquotedblleft{}The calf of the leg\textquotedblright{}.) Say it once more.
\end{tcolorbox}
\section[\texorpdfstring{\gu{બગલ} (bagal)}{બગલ (bagal)}]{\gu{બગલ} (bagal) — an armpit}

\label{lesson:GU-C473-bagal}



\begin{quote}
Before the new one: say the Gujarati for the navel, then the Gujarati for the calf of the leg.
\end{quote}

\subsection*{You'll want to know: \gu{બગલ}}
\textbf{\gu{બગલ}} — \emph{bagal} — \textquotedblleft{}an armpit\textquotedblright{}.

\begin{quote}
\gu{નર્સે} \gu{બગલમાં} \gu{તાવ} \gu{માપ્યો}.
\end{quote}

\begin{quote}
\emph{narse bagalmā̃ tāv māpyo.}
\end{quote}

\begin{quote}
The nurse took my temperature under my arm.
\end{quote}

A feminine noun.

\begin{cousinweb}
\textbf{\gu{બગલ}} (\emph{bagal}) is usually traced to Persian \emph{baġal}, \textquotedblleft{}armpit, side\textquotedblright{}.
\end{cousinweb}

\begin{grammarlens}[title={What you've built}]
One more word for describing what you lived and felt.
\end{grammarlens}

\subsection*{Your turn}
\begin{itemize}
\raggedright
  \item \emph{Say it:} \emph{bagal}
  \item \emph{Say it:} \emph{bagal}, once more
  \item \emph{Say it:} \emph{narse bagalmā̃ tāv māpyo.} — The nurse took my temperature under my arm.
  \item \emph{Recall:} say the Gujarati for the navel, then the Gujarati for the calf of the leg, then say \emph{bagal} again
\end{itemize}

\begin{tcolorbox}[breakable,colback=teal!4,colframe=teal!35!black,title={Before you move on}]
What does \gu{બગલ} mean? (\textquotedblleft{}An armpit\textquotedblright{}.) Say it once more.
\end{tcolorbox}
